\begin{abstract}
More than $1.5$ billion people live with hearing loss, and when several people
talk at once, a hearing aid amplifies every voice alike. Auditory attention
decoding (AAD) aims to solve this: it reads a listener's brain signals, finds the
talker they are attending to, and lets the device amplify that talker. A real
device will not see the brain alone. Glasses and earbuds can also track the eyes,
the head and the scene, and these signals help, because people tend to look at
the person they listen to. But they show where a listener is looking, not whom
they are listening to. A decoder that leans on them works when the two agree and
fails when they differ --- when a listener cannot move their eyes or head, or
listens to someone they are not looking at --- which is exactly when a brain-based
device is needed. A multimodal decoder must therefore use these signals without
letting them replace the EEG, and accuracy alone cannot show whether it does. We
present \merit{}, a decoder and an evaluation method built around one question:
how much of the accuracy does each input earn? \merit{} shuffles each input
stream across test examples, measures how much accuracy drops, and splits the
result exactly across the streams with Shapley values. It also blocks two
shortcuts by design: the candidate audio is checked before training so that it
cannot reveal the answer, and the scoring function cannot rise above chance if it
ignores the brain signal. On a public dataset with EEG, eye tracking, head motion
and scene video, \merit{} improves on the published benchmark, reaching
\R{v2hinge2/H4_nohinge_nodrop/within/acc2} four-way accuracy within listeners and
\R{v2hinge2loso/H4_nohinge_nodrop/loso/acc2} on unseen listeners, where chance is
$0.25$. The measurement also shows what accuracy hides: adding the behavioural
streams raises accuracy but almost halves the brain's share of it. On two other
public datasets, the same tests find a label imbalance in a widely used benchmark
and flag results that accuracy alone would report as brain decoding.
\end{abstract}
